\chapter{Implementacja}
\label{chap:implementation}
\section{Implementacja wyglądu i funkcjonalności frontendu}
W tej sekcji omówione zostaną kluczowe aspekty implementacji frontendu aplikacji, w tym architektura, struktura UI i komponentów, warstwa wizualna, obsługa interakcji użytkownika, integracja z geolokalizacją, implementacja trybów gry, komunikacja z backendem oraz mechanizmy PWA.

\subsection{Architektura frontendu (Karina Wołoszyn)}

\subsection{Struktura UI i komponentów (Karina Wołoszyn)}

\subsection{Warstwa wizualna (Karina Wołoszyn)}

\subsection{Obsługa interakcji użytkownika (Karina Wołoszyn)}

\subsection{Integracja z geolokalizacją (Filip Jezierski)}

\subsection{Implementacja trybów gry (Filip Jezierski)}

\subsection{Komunikacja z backendem (Filip Jezierski)}

\subsection{Mechanizmy PWA (Filip Jezierski)}

% \subsection{Problemy napotkane  przy implementacji frontendu}


\section{Implementacja backendu i logiki aplikacji}

W tej sekcji omówione zostaną kluczowe aspekty implementacji backendu aplikacji, rozgrywki  mechanizmy autoryzacji i uwierzytelniania użytkowników, zarządzanie danymi oraz obsługa błędów, a także problemy, które napotkano podczas implementacji.

\subsection{Implementacja rozgrywki}

Przedstawienie sposobu w jaki została zaimplementowana rozgrywka z \ref{fig:sequence-diagram} w aplikacji. 

Proces rozpoczęcia rozgrywki inicjuje wywołanie kontrolera \texttt{api/game/start}, który odpowiada za utworzenie nowej \textit{Sesji Rozgrywki} (\ref{tab:game-sessions}).
W ramach tego etapu backend generuje odpowiednią liczbę rund (\ref), przypisuje im losowe miejsca do zgadnięcia, nadaje sesji status \texttt{InProgress}, a następnie zwraca unikalny identyfikator gry w postaci GUID.
W przypadku użytkownika niezalogowanego aplikacja dodatkowo generuje i zwraca tymczasowy token, który umożliwia jego uwierzytelnienie w kolejnych żądaniach.

Aby pobrać dane dotyczące danej rundy, frontend wywołuje endpoint
\texttt{api/game/{gameId}/round/{roundNumber}},
gdzie \texttt{gameId} jest identyfikatorem sesji, a \texttt{roundNumber} – numerem aktualnej rundy. Backend zwraca ścieżkę do obrazu przedstawiającego miejsce, które należy odgadnąć. Obraz może być przechowywany lokalnie na serwerze lub stanowić odnośnik do zewnętrznego źródła.


Odpowiedź użytkownika jest obsługiwana przez endpoint
\texttt{api/game/{gameId}/round/{roundNumber}/check}.
Aplikacja weryfikuje lokalizację podaną przez użytkownika, oblicza odległość pomiędzy faktycznym miejscem a wskazanym punktem oraz na tej podstawie przyznaje punkty za daną rundę.

Po zakończeniu wszystkich rund gracz wywołuje endpoint
\texttt{api/game/{gameId}/finish},
który finalizuje rozgrywkę. Backend podsumowuje wyniki, aktualizuje status sesji na \texttt{Completed} i zwraca końcowy rezultat użytkownikowi.
Dla użytkowników zalogowanych rezultat jest trwale zapisywany w bazie, natomiast dla użytkowników niezalogowanych sesja jest automatycznie usuwana po upływie jednej godziny.


\subsection{Autoryzacja i uwierzytelnianie użytkowników}


\textbf{Generowanie i struktura tokenu JWT służącego do autoryzacji użytkowników}

W celu autoryzacji wykorzystany został JSON Web Token. Jest to standard wymiany informacji między dwoma stronami w formie obiektu JSON, który jest bezpiecznie podpisany cyfrowo. 

Składa się on z trzech niezależnych części zakodowanych w formacie Base64: nagłówka (header), ładunku (payload) oraz podpisu (signature). Nagłówek zawiera informacje o algorytmie podpisu i typie tokenu, ładunek zawiera dane użytkownika i inne informacje, a podpis służy do weryfikacji integralności tokenu i autentyczności nadawcy.

Przykładowy token JWT:
\begin{lstlisting}
eyJhbGciOiJIUzI1NiIsInR5cCI6IkpXVCJ9.eyJzdWIiOiIxMjM0NTY3ODkwIiwibmFtZSI
6IkpvaG4gRG9lIiwiYWRtaW4iOnRydWUsImlhdCI6MTUxNjIzOTAyMn0.KMUFsIDTnFmyG3n
MiGM6H9FNFUROf3wh7SmqJp-QV30
\end{lstlisting}

Nagłówek przedstawia zakodowane informacje o algorytmie podpisu oraz typie tokenu (w tym przypadku HS256):

\begin{lstlisting}
{
  "alg": "HS256",
  "typ": "JWT"
}
\end{lstlisting}

Payload zawiera zakodowane informacje, które chcemy przekazać, jak np. identyfikator użytkownika, nazwa użytkownika, role itp.:

\begin{lstlisting}
{
  "ID": "23141",
  "name": "Jan Kowalski",
  "admin": true,
  "iat": 1516239022
}
\end{lstlisting}

Podpis jest generowany przez zakodowanie nagłówka i payloadu oraz podpisanie ich kluczem prywatnym przy użyciu algorytmu określonego w nagłówku.

W środowisku .NET mechanizm tokenów JWT jest wspierany przez bibliotekę \texttt{Microsoft.AspNetCore.Authentication.JwtBearer}, która umożliwia łatwą integrację z mechanizmami uwierzytelniania i autoryzacji w aplikacjach ASP.NET Core. Dzięki temu jesteśmy w stanie generować w efektywny sposób tokeny, które wykorzystywane są do autoryzacji. 

\begin{lstlisting}[language={[Sharp]C}, caption={Generowanie tokenu JWT w warstwie backendu}, label={lst:jwt-generation}]
var claims = new List<Claim>
{
    new(ClaimTypes.NameIdentifier, user.Id.ToString()),
    new(ClaimTypes.Email, user.Email),
    new(ClaimTypes.Role, user.Role),
    new(ClaimTypes.Name, user.Nickname),
    new("token_type", "user")
};

// Klucz podpisujący token (HMAC SHA-256)
var key = new SymmetricSecurityKey(
    Encoding.UTF8.GetBytes(_authenticationSettings.JwtKey));

// Poświadczenia podpisu
var credentials = new SigningCredentials(
    key, SecurityAlgorithms.HmacSha256);

// Czas wygaśnięcia tokenu
var expires = DateTime.Now.AddMinutes(
    _authenticationSettings.JwtExpireAccount);

// Utworzenie obiektu tokenu JWT
var token = new JwtSecurityToken(
    issuer: _authenticationSettings.JwtIssuer,
    audience: _authenticationSettings.JwtIssuer,
    claims: claims,
    expires: expires,
    signingCredentials: credentials
);
\end{lstlisting}

\textbf{Identyfikowanie użytkowników i autoryzacja dostępu do zasobów}

W aplikacji wykorzystano mechanizm autoryzacji oparty na atrybutach \texttt{[Authorize]} dostępnych w ASP.NET Core. Atrybut ten pozwala na zabezpieczenie kontrolerów lub konkretnych akcji, wymagając od użytkownika posiadania ważnego tokenu JWT oraz odpowiednich uprawnień (ról) do dostępu do zasobów. 
Na etapie przetwarzania żądania, middleware uwierzytelniające sprawdza obecność i ważność tokenu JWT w nagłówku autoryzacji. Jeśli token jest poprawny, użytkownik zostaje zidentyfikowany, a jego rola jest sprawdzana w kontekście żądanej akcji.


\textbf{Hashowanie haseł użytkowników i uwierzytelnianie}

Do bezpiecznego generowania i przechowywania haseł użytkowników wykorzystano .NET Core Identity, który implementuje algorytm HMAC-SHA256. Algorytm ten jest odporny na ataki brute-force dzięki zastosowaniu kosztownego obliczeniowo procesu hashowania oraz losowego saltingu. W procesie rejestracji hasło użytkownika jest hashowane przed zapisem do bazy danych, natomiast podczas logowania następuje weryfikacja hasła poprzez porównanie jego hasha z wartością przechowywaną w bazie.

\begin{lstlisting}[language={[Sharp]C}, caption={Hashowanie i weryfikacja hasła użytkownika}]
// Hashowanie hasła podczas rejestracji
var passwordHash = passwordHasher.HashPassword(
    newUser, 
    registerAccountCommand.Password
);

// Weryfikacja hasła podczas logowania
var result = passwordHasher.VerifyHashedPassword(
    user, 
    user.PasswordHash, 
    loginCommand.Password
);
\end{lstlisting}

\subsection{Walidacja danych}
    W celu zapewnienia poprawności danych wejściowych w aplikacji zastosowano bibliotekę FluentValidation. Umożliwia ona definiowanie reguł walidacyjnych w sposób deklaratywny, czytelny i łatwy w utrzymaniu. Przykładowo, podczas rejestracji użytkownika można zdefiniować reguły dotyczące minimalnej i maksymalnej długości hasła, prawidłowego formatu adresu e-mail oraz unikalności nazwy użytkownika w systemie. Dzięki temu możliwe jest centralne zarządzanie walidacją oraz zapewnienie spójności danych w całej aplikacji. 
    
    Biblioteka FluentValidation jest w pełni zintegrowana z mechanizmem wiązania modeli (model binding) w ASP.NET Core, co umożliwia automatyczne wywoływanie walidacji podczas przetwarzania żądań HTTP. W przypadku wykrycia błędów walidacji framework automatycznie zwraca odpowiedź HTTP 400 (Bad Request) wraz ze szczegółowym opisem napotkanych problemów, co znacząco ułatwia debugowanie i poprawia komunikację z klientem API.

    Niektóre reguły walidacyjne wymagają dostępu do warstwy persystencji, na przykład w celu sprawdzenia, czy w bazie danych nie istnieje już użytkownik o podanym adresie e-mail lub nazwie. W takich przypadkach do konstruktora klasy walidatora wstrzykiwane są odpowiednie repozytoria poprzez mechanizm wstrzykiwania zależności (Dependency Injection). 

\begin{lstlisting}[language={[Sharp]C}, caption={Przykład walidacji danych użytkownika podczas rejestracji}, label={lst:fluent-validation}]
public class RegisterUserDtoValidator : AbstractValidator<RegisterUserDto>
{
    private readonly IAccountRepository _accountRepository;

    public RegisterUserDtoValidator(IAccountRepository accountRepository)
    {
        _accountRepository = accountRepository;
        RuleFor(u => u.Email)
            .NotEmpty()
            .EmailAddress();

        RuleFor(u => u.Nickname)
            .NotEmpty()
            .MinimumLength(3)
            .MaximumLength(25);

        RuleFor(u => u.Password)
            .NotEmpty()
            .MinimumLength(8)
            .WithMessage("Minimum length of password is 8");

        RuleFor(x => x.ConfirmPassword)
            .Equal(e => e.Password)
            .WithMessage("Passwords are not the same ");
    }
}
\end{lstlisting}


Dzięki takiemu podejściu walidacja danych jest centralnie zarządzana i łatwo rozszerzalna, co przyczynia się do poprawy jakości oraz bezpieczeństwa aplikacji. Dodatkowo każda reguła walidacyjna jest umieszczona w osobnej klasie, co pozwala przestrzegać zasady pojedynczej odpowiedzialności oraz zachować czytelność i modularność kodu. Umożliwia to również łatwe testowanie poszczególnych walidatorów w izolacji od pozostałych komponentów systemu.


% potencjalnie napisze moze cos o mediator i flunet validation tutaj

\subsection{Obsługa błędów w aplikacji}
    W aplikacji do obsługi błędów wykorzystano mechanizm dostępyn w ASP .NET Core, polegający na tworzeniu globalnego middleware odpowiedzialnego za przechwytywanie wyjątków. Middleware jest ten w stanie obsłużyć wyjątki rzucane w dowolnym miejscu w aplikacji i zwrócić odpowiednią odpowiedż z kodem statusu HTTP oraz komunikatem błędu do klienta. Dzieki temu użytkownik otrzymuje spójne i czytelne informacje o błędach, bez widocznego stosu wywołan i komunikatów serwera.  Dzięki temu podejściu możemy łatwo zarządzać obsługą błędów w jednym miejscu i zapewnić lepsze doświadczenie użytkownika.

    Dodatkowo dla każdej encji w warstwie domeny zdefiniowano oddzielny plik z wyjątkami, które mogą wystapić podczas operacji na tych encjach. Dzięki temu łatwiej jest zarządzać różnymi typami błędów, charakterystycznymi dla danej encji.




\subsection{Warstwa persystencji - Entity Framework Core}
    
    -Rola orm w aplikacji
    -DbContext i mapowanie encji na tabele bazy danych
    -Migracje bazy danych
    -Repozytoria w aplikacji

\subsection{Problemy napotkane przy implementacji backendu}
    -Eager loading vs lazy loading w EF Core
    -Konflikty migracji bazy danych
    -Obsługa błędów w aplikacji ASP .NET Core
    -Debuggowanie aplikacji w kontenerze Docker

\section{Baza danych i zarządzanie danymi}



\subsection{Postgres i adminer}


\subsection{Przechowywanie plików }




\section{Devops i wdrożenie aplikacji}

W tej sekcji opisano proces przygotowania, uruchamiania oraz wdrażania aplikacji w środowiskach deweloperskim i produkcyjnym z wykorzystaniem konteneryzacji Docker, plików konfiguracyjnych oraz reverse proxy Nginx. Przedstawiono również decyzje projektowe dotyczące zarządzania certyfikatami oraz publikacji obrazów w zewnętrznym rejestrze (Docker Hub).

\subsection{Środowisko deweloperskie}

Środowisko deweloperskie zostało zaprojektowane tak, aby umożliwić szybkie iteracje oraz łatwe debugowanie. Kluczowe założenia:
\begin{itemize}
    \item Oddzielne kontenery dla frontendu (React/TypeScript) i backendu (ASP.NET Core) uruchamiane w trybie \texttt{watch/hot reload}.
    \item Kontener bazy danych PostgreSQL oraz (opcjonalnie) narzędzie \texttt{Adminer} do inspekcji danych.
    \item Użycie pliku \texttt{docker-compose.dev.yaml} do orkiestracji wielokontenerowego środowiska.
\end{itemize}

Konfiguracja usług w pliku compose dla środowiska deweloperskiego może zawierać sekcje z wolumenami mapowanymi na katalogi źródłowe, aby kod był automatycznie odświeżany. Backend korzysta z \texttt{dotnet watch run}, frontend z \texttt{npm start}. Daje to krótki cykl: edycja $\rightarrow$ zapis $\rightarrow$ natychmiastowy efekt. Eliminuje to długie czasy rekompilacji całej aplikacji.

\paragraph{Debugowanie w kontenerze} Dla backendu możliwe jest podłączenie zdalnego debuggera (np. VS Code / Rider) przez port udostępniony w kontenerze. W środowisku deweloperskim certyfikaty HTTPS nie są wymagane, ponieważ komunikacja odbywa się po \texttt{http://localhost}.

\subsection{Środowisko produkcyjne}

W środowisku produkcyjnym nacisk położono na izolację usług, bezpieczeństwo i skalowalność. Kluczowe elementy:
\begin{itemize}
    \item Użycie pliku \texttt{docker-compose.prod.yaml} z obrazami już zbudowanymi (bez mountowania kodu źródłowego).
    \item Wprowadzenie reverse proxy Nginx jako pojedynczego punktu wejścia (terminacja TLS, routing do usług wewnętrznych).
    \item Minimalizacja rozmiaru obrazów przez wykorzystanie wieloetapowych buildów (multi-stage) w Dockerfile backendu i frontendu.
    \item Brak narzędzi administracyjnych (np. Adminer) w docelowym środowisku - uruchamiane jedynie tymczasowo przy potrzebie migracji lub analizy.
\end{itemize}

\subsection{Pliki docker-compose i ich role}

W projekcie zastosowano kilka wariantów plików \texttt{docker-compose}. Mają one strukturę kaskadową:
\begin{description}
    \item[\texttt{docker-compose.infra.yaml}] Definicje infrastrukturalne (baza danych, narzędzia pomocnicze). Pozwala odseparować warstwę aplikacyjną od zasobów zewnętrznych.
    \item[\texttt{docker-compose.dev.yaml}] Konfiguracja zawierająca część frontendową oraz infrastrukturalną, umożliwia uruchomienie backendu poza kontenerem. 
    \item[\texttt{docker-compose.yaml}] Plik zawierający wszystkie powyższe rzeczy, gotowy do uruchomienia przez członków zespołu bez zbędnej konfiguracji.
    \item[\texttt{docker-compose.prod.yaml}] Uporządkowane obrazy z tagami wersji, brak wolumenów źródłowych, zmienne środowiskowe wczytywane z plików \texttt{.env}, integracja z Nginx.
\end{description}

co pozwala na nadpisywanie sekcji (\textit{override}) dla danego środowiska.

\subsection{Decyzja o użyciu reverse proxy Nginx}

Pierwotnie rozważano generowanie i instalację certyfikatów HTTPS w każdej części aplikacji osobno (frontend i backend). Takie rozwiązanie zwiększało złożoność utrzymania (odnawianie, dystrybucja, synchronizacja dat ważności). Zamiast tego wprowadzono warstwę \texttt{Nginx} jako reverse proxy z terminacją TLS:
\begin{itemize}
    \item Jeden certyfikat (np. Let's Encrypt) dla domeny głównej.
    \item Uproszczona konfiguracja \texttt{appsettings.json} - backend może działać na HTTP wewnątrz sieci Docker.
    \item Brak potrzeby modyfikowania konfiguracji frontendu przy odnowieniu certyfikatu - zmiana dotyczy wyłącznie proxy.
    \item Możliwość dodania reguł bezpieczeństwa (rate limiting, nagłówki CSP, HSTS) w jednym miejscu.
\end{itemize}

\paragraph{Korzyści architektoniczne} Centralizacja TLS redukuje liczbę punktów potencjalnych błędów i przyspiesza automatyzację (skrypt odnawiający certyfikat uruchamiany cyklicznie w kontenerze \texttt{certbot}).

\subsection{Certyfikaty i ich zarządzanie}

W katalogu \texttt{https/} znajdował się lokalny certyfikat deweloperski \texttt{dev-cert.pfx} używany w fazie wczesnego rozwoju do testów lokalnych. Po wdrożeniu Nginx rola tego pliku została ograniczona wyłącznie do scenariuszy developerskich. W produkcji certyfikaty są pobierane automatycznie (np. przez \texttt{certbot}) i przechowywane w zamontowanym wolumenie na hosta, dzięki czemu ich odnowienie nie wymaga przebudowy obrazów.

\subsection{Budowa i publikacja obrazów Docker}

Proces publikacji obrazu na \texttt{Docker Hub} obejmuje kroki:
\begin{enumerate}
    \item Budowa obrazu (np. \texttt{docker build -t użytkownik/nazwa:tag .}).
    \item Logowanie do rejestru (\texttt{docker login}).
    \item Wysłanie obrazu (\texttt{docker push użytkownik/nazwa:tag}).
    \item Opcjonalne użycie wersjonowania semantycznego (\texttt{1.2.0}, \texttt{latest}).
\end{enumerate}

Multi-stage build pozwala oddzielić etap kompilacji (cięższy obraz z SDK) od etapu wykonawczego (lżejszy runtime). Zmniejsza to czas pobierania i powierzchnię ataku.

\paragraph{Tagowanie i automatyzacja} Przy każdej zmianie głównej funkcjonalności można generować nowy tag. W przyszłości możliwe jest dodanie pipeline CI (GitHub Actions) automatyzującej budowę i publikację na podstawie tagów w repozytorium Git.

\subsection{Konfiguracja aplikacji}

Konfiguracja jest rozproszona pomiędzy kilka plików i mechanizmów:
\begin{itemize}
    \item Pliki \texttt{.env} Zawierają wrażliwe lub zmienne środowiskowe (np. connection string do PostgreSQL, sekrety JWT). W plikach compose referencje przez \texttt{\$\{NAZWA\_ZMIENNEJ\}}.
    \item \texttt{appsettings.json}, \texttt{appsettings.Development.json}, \texttt{appsettings.Production.json} Warstwowe źródło konfiguracji backendu: logowanie, połączenie do bazy, parametry JWT (issuer, audience, czas życia), ustawienia CORS.
    \item \texttt{launchSettings.json} Używany lokalnie (poza Docker) do definiowania profili uruchomieniowych i portów (HTTPS/HTTP). W produkcji zastąpione przez konfigurację kontenera + reverse proxy.
    \item Frontend \texttt{Constants.ts} Zawiera stałe takie jak URL API, nazwy eventów, parametry mapy. W produkcji wartości mogą być nadpisane przez zmienne środowiska przy budowie (np. \texttt{VITE\_API\_URL}).
    \item Zmienne środowiskowe kontenerów Nadawane w plikach compose lub w CI/CD 
    (np. \texttt{ASPNETCORE\_ENVIRONMENT=Production}).
\end{itemize}

Mechanizm konfiguracji w ASP.NET Core scalający źródła (pliki, zmienne środowiskowe, sekrety użytkownika) umożliwia elastyczne przełączanie środowisk bez modyfikacji kodu.

\paragraph{Bezpieczeństwo konfiguracji} Sekrety (klucz do podpisu JWT, hasło do bazy) nie powinny być wersjonowane --- w środowisku deweloperskim mogą być przechowywane w pliku \texttt{.env.local} dodanym do \texttt{.gitignore}. W produkcji zalecane jest ich wstrzykiwanie poprzez menedżera sekretów (np. HashiCorp Vault / Docker Swarm secrets / Kubernetes secrets).

\subsection{Przepływ uruchomienia aplikacji}

Cały proces startu w produkcji można streścić następująco:
\begin{enumerate}
    \item Nginx startuje i ładuje certyfikat TLS.
    \item Backend uruchamia się w trybie HTTP wewnątrz sieci Docker, migracje bazy wykonywane automatycznie na starcie (opcjonalnie sterowane flagą).
    \item Frontend serwowany jako statyczne pliki (np. z \texttt{nginx}) lub z dedykowanego kontenera serwera plików.
    \item Nginx przekazuje żądania \texttt{/api} do backendu, a pozostałe ścieżki obsługuje przez serwowanie SPA.
\end{enumerate}

\subsection{Podsumowanie sekcji DevOps}

Zastosowana architektura DevOps upraszcza cykl życia aplikacji: deweloperzy otrzymują szybkie iteracje dzięki mountowaniu kodu, produkcja korzysta z odchudzonych i bezpiecznych obrazów oraz centralnego zarządzania TLS. Decyzja o użyciu reverse proxy Nginx z jednym certyfikatem poprawiła bezpieczeństwo i zmniejszyła nakład operacyjny, tworząc solidną bazę pod dalsze skalowanie i obserwowalność.



\section{Integracja komponentów systemu}

